\begin{exo}{}
	$f$ est une fonction définie sur $\intff{-2}{2}$ dont on donne le tableau de valeurs suivant.


	\begin{center}
		\renewcommand{\arraystretch}{1.5}
		\small
		\begin{tabular}{|c|c|c|c|c|c|}\hline
			\cellcolor{gray!30}$x$    & $-2$  & $-1,5$ & $-1$  & $-0,5$ & $0$ \\ \hline
			\cellcolor{gray!30}$f(x)$ & $-1$  & $-3$   & $-1$  & $3$    & $2$ \\ \hline\hline
			\cellcolor{gray!30}$x$    & $0,5$ & $1$    & $1,5$ & $2$    &     \\ \hline
			\cellcolor{gray!30}$f(x)$ & $0$   & $2$    & $1$   & $3$    &     \\ \hline
		\end{tabular}
	\end{center}

	Quatre élèves proposent ci-dessous une représentation graphique de $f$.

	\begin{center}
		\begin{tikzpicture}[scale=0.75, >=Stealth]
			% Configuration commune des points
			\def\pts{ (-2, -1), (-1.5, -3), (-1, -1), (-0.5, 3), (0, 2), (0.5, 0), (1, 2), (1.5, 1), (2, 3) }
			\tikzset{
				grid style/.style={gray!40, thin},
				axis style/.style={thick, ->},
				pt style/.style={fill=purple!80!black, draw=purple!80!black, circle, inner sep=0pt, minimum size=4pt}
			}
			\newcommand{\drawGridAndAxes}{
				\draw[grid style, xstep=1, ystep=1] (-2.5, -3.5) grid (2.5, 3.5);
				\draw[axis style] (-2.5, 0) -- (2.5, 0);
				\draw[axis style] (0, -3.5) -- (0, 3.5);
				\foreach \x in {-2,-1,1,2} \draw (\x, 0.1) -- (\x, -0.1);
				\foreach \y in {-3,-2,-1,1,2,3} \draw (0.1, \y) -- (-0.1, \y);
				\node[below left] at (0, 0) {$O$};
				\node[below] at (1, 0) {$I$};
				\node[left] at (0, 1) {$J$};
			}
			\begin{scope}[shift={(0, 7.5)}]
				\drawGridAndAxes
				\foreach \coord in \pts { \node[pt style] at \coord {}; }
			\end{scope}

			% --- SCOPE 2 : Haut-Droite (Ligne brisée) ---
			\begin{scope}[shift={(5.5, 7.5)}]
				\drawGridAndAxes
				\draw[thick, purple!80!black]
				(-2, -1) -- (-1.5, -3) -- (-1, -1) -- (-0.5, 3) -- (0, 2) -- (0.5, 0) -- (1.5, 1) -- (1, 2) -- (2, 3);
				\foreach \coord in \pts { \node[pt style] at \coord {}; }
			\end{scope}

			% --- SCOPE 3 : Bas-Gauche (Courbe restreinte à l'intervalle) ---
			\begin{scope}[shift={(0, 0)}]
				\drawGridAndAxes
				\draw[thick, purple!80!black, smooth] plot coordinates { (-2, -1) (-1.5, -3) (-1, -1) (-0.5, 3) (0, 2) (0.5, 0) (1, 2) (1.5, 1) (2, 3) };
				\foreach \coord in \pts { \node[pt style] at \coord {}; }
			\end{scope}

			% --- SCOPE 4 : Bas-Droite (Courbe prolongée aux extrémités) ---
			\begin{scope}[shift={(5.5, 0)}]
				\drawGridAndAxes
				\clip (-3.5, -3.5) rectangle (2.5, 3.5);
				\draw[thick, purple!80!black, smooth] plot coordinates { (-2.5, 3) (-2, -1) (-1.5, -3) (-1, -1) (-0.5, 3) (0, 2) (0.5, 0) (1, 2) (1.5, 1) (2, 3) (2.2, 3.5) };
				\foreach \coord in \pts {
					\node[pt style] at \coord {};
				}
			\end{scope}

		\end{tikzpicture}
	\end{center}


	Laquelle de ces courbes semble la plus satisfaisante ? Justifier.
\end{exo}
